\chapter{Orbital portrait (menu 15)}
\label{ch:menu15}
\menulabel{15}

\section{Purpose}

The deterministic part of RLEASE's 26-dimensional orbital descriptor -- the part
that needs no model and no training, which is what this project keeps (the
reading notes record that the source's predictor network and threshold policy
are explicitly not adopted).  The panel reads, per orbital: the occupation and
the canonical energy; the dominant centre with its L\"owdin share; the
angular-momentum composition; the \textbf{bonding label}; and two AO-centre
estimates (the spatial extent and the charge-centroid displacement).

\section{What you need}

One \code{orca\_2json} export (the MO coefficients, the occupations, the overlap
and the AO labels with the atom coordinates).  The window defaults to the
orbitals with fractional occupations, i.e. a correlated export's active space.

\section{Walkthrough}

\lstinputlisting[style=fbkcode, firstline=19, lastline=39,
  title={The menu-15 example session (the N\textsubscript{2} CAS(6,6) export,
  inferred active window)}]
  {examples/expected/m15_portrait.txt}

\section{What you get}

The table on screen and a report file (\file{<export>.portrait.fbk.md}) with the
criteria and the citations.

\section{How to read the numbers}

\begin{readbox}
\begin{itemize}
  \item \textbf{The bonding label is the source's own criterion.}  Accumulate
    the atom-pair population $S_{AB} = \sum_{\mu\in A}\sum_{\nu\in B} c_\mu
    S_{\mu\nu} c_\nu$ over atom pairs within $6$~\AA; a single atom carrying
    more than $95\%$ of the L\"owdin population makes the orbital non-bonding;
    otherwise the sign of the cross terms decides bonding from antibonding.
    Two measured additions travel with it: a cross term within $0.01$ counts as
    non-bonding (a symmetric $1s$ core measures $3\times10^{-4}$, and a bare
    sign test would call such a core ``bonding''), and a row marked \code{!} is
    \emph{cancellation-heavy} -- the absolute sum of its atom-pair blocks
    exceeds $5$ -- so its cross magnitude is not a bond order even though its
    sign still reads.
  \item \textbf{The shares are L\"owdin atomic populations}: non-negative and
    summing to one for every orbital.  The raw block sums are \emph{not}
    interchangeable with them: the same node-heavy virtual that carries the
    cancellation mark has block sums of $40$.
  \item \textbf{The extent and the centroid displacement are AO-centre
    estimates.}  The source's exact $\langle r^2\rangle$ and its diagonal
    integrals need matrix elements an orbital export does not carry; an FCIDUMP
    has the latter when they are wanted (menu~\menu{12}'s route).
  \item \textbf{The composition is the weighted form of the source's binary
    shell encoding.}  The source deliberately encodes shells as present/absent
    to stay insensitive to the basis size; the reading notes record that this
    project's shell-classification checks need the weighted version, so the
    panel prints weights.
  \item \textbf{The panel is evidence, not the choice.}  The source's own
    benchmark is the warning here: a larger active space is not automatically a
    better one (on CH\textsubscript{4} the reference method picked 17 orbitals
    and came out with twice the error of an 8-orbital space), so a risk reading
    that treats ``active space smaller than expected'' as a defect on its own
    would be wrong -- the project's diagnosis rules do not do that.
\end{itemize}
\end{readbox}

\section{Worked example}

The session above runs on the shipped N\textsubscript{2} export and reads its
chemistry back: the two $1s$-derived cores come out non-bonding, the $\pi$ pair
bonding, the $\pi^*$ pair antibonding, and exactly one node-heavy virtual
carries the cancellation mark.  That is the panel's purpose in one line: before
choosing a window, see what the orbitals are.

\section{Caveats, refusals and related sections}

\begin{refusalbox}
\begin{itemize}
  \item An export without an overlap matrix or without AO labels is refused.
  \item A window outside the export's orbital range, or an empty one, is
    refused.
  \item The panel is a characterisation, not a selection: it neither ranks
    orbitals by importance nor recommends a space (menu~\menu{14} sizes one,
    menu~\menu{13} compares two).
\end{itemize}
\end{refusalbox}

Related: \menu{14} (the target-shell projection that sizes a space),
\menu{13} (the space comparison, whose report now carries the Jaccard index of
the two windows -- the process indicator the descriptor-panel source's
evaluation protocol uses next to the energy error), and
Chapter~\ref{ch:criteria} for the threshold table.
